\section{Odd cyclotomic orbit sizes for induction columns}
\label{sec:odd-column-orbits}
The relative twist traces can be strengthened by retaining their orbit sizes.
This refinement is a necessary condition for a pseudo-unitary categorification;
a feasible system is not a construction of a fusion category.

\begin{proposition}\label{prop:odd-column-orbits}
Let $\mathcal C$ be pseudo-unitary with commutative Grothendieck ring, and let
$K=\mathbb Q(\zeta_F)$ contain its complete character table and its Frobenius--Perron dimensions.
Assume that $F$ is odd and that every prime divisor of the Frobenius--Schur
exponent divides $F$.  Fix an induction restriction column $v$, and let $m_v$
be the number of simple objects of $\mathcal Z(\mathcal C)$ having this column.
For a twist of order $n$, put
\[
 e(n)=[\mathbb Q(\zeta_{\operatorname{lcm}(F,n)}):K]
      =\prod_{p\mid F}p^{\max(0,v_p(n)-v_p(F))}.
\]
The multiset of twists over $v$ is a disjoint union of full relative Galois
orbits.  In particular an orbit of twists of order $n$ contributes $e(n)$
objects, so $e(n)\le m_v$.  Its contribution to the first $h$ twist moments is
\[
 e(n)\left(\frac{\operatorname{Tr}_{K(\zeta_n)/K}(\zeta_n^a)}{e(n)},\ldots,
 \frac{\operatorname{Tr}_{K(\zeta_n)/K}(\zeta_n^{ah})}{e(n)}\right).
\]
\end{proposition}

\begin{proof}
Write $D=\operatorname{FPdim}(\mathcal C)$ and let $A_\chi$ denote the
multiplicity-one simple summand of $I(\mathbf1)$ associated to a character
$\chi$, of formal codegree $f_\chi$.  With normalized $S$-matrix,
\[
 S_{A_\chi,Z}=\frac{\chi(FZ)}{f_\chi},\qquad
 S_{\mathbf1,Z}=\frac{\operatorname{FPdim}(FZ)}{D}.
\]
Thus all these entries belong to $K$.  A modular Galois automorphism fixing
$K$ fixes the unit row and every $A_\chi$ row: its possible sign is positive,
as seen in the positive unit column, and distinct simple objects have
 distinct $S$-matrix rows.  It also fixes the entries in these rows when
permuting any column $Z$.  Invertibility of the ordinary character table
therefore shows that $FZ$ is unchanged.  This proves invariance of each
restriction-column class.

For the canonical ribbon structure on a Drinfeld center, the modular Galois
permutation transforms a twist by $\theta\mapsto\sigma^2(\theta)$, as in the modular Galois symmetry of \cite{BNRW}; there is
no anomaly factor.  The group
$\operatorname{Gal}(K(\zeta_n)/K)$ has the displayed order $e(n)$, which is
odd, and the square map is an automorphism of this abelian group.  Hence its
square action on primitive $n$th roots has the same orbits as its ordinary
Galois action.  Each such orbit has $e(n)$ elements, with equal
multiplicities in the fixed restriction-column class.  Summing the powers
of these roots proves the formula.
\end{proof}

Here is the finite form used in the supplementary implementation.  If $G$
is the residual induction Gram matrix after the low columns have been
removed, a nonzero high column $v$ has multiplicity at most
\[
 B_v=\min_{i:v_i\ne0}\left\lfloor\frac{G_{ii}}{v_i^2}\right\rfloor.
\]
Consequently only orders $n$ with $e(n)\le B_v$ are relevant.  They are
exhaustively enumerated by
$0\le v_p(n)\le v_p(F)+\lfloor\log_p B_v\rfloor$.
For each normalized trace type, the count must belong to the additive
semigroup generated by the corresponding orbit sizes.  Removing a size
which is a positive multiple of a smaller size does not change this
semigroup.  The code uses a separate nonnegative integer variable for each
retained size, multiplying the corresponding moment column by that size.
No divisibility assertion is made when the relative Galois group has even
order.

\subsection{The conductor-fifteen multiplicity-four frontier}
For $\RR{4}{629}$, the formal-codegree polynomial is
\[
 (t-5)^2(t^4-360t^3+5850t^2-30375t+50625).
\]
The unrestricted signed enumeration has exactly $32$ tables in one sign
orbit, all gauge-equivalent to the original table.  The exact dimension
character test leaves only the Frobenius--Perron Galois orbit.  Thus every
categorification is spherical and has a pseudo-unitary Galois conjugate;
the following necessary-data reduction therefore applies to general
categorifications as well.

There are seven low induction frames.  Joint-spectrum integrality eliminates
six.  The remaining frame has two possible character labellings and $85$
admissible high restriction vectors.  Applying
Proposition~\ref{prop:round4-affine-profile-enumeration} to its Gram and
Galois equations gives exactly $784$ nonnegative integral column-multiplicity
profiles.  Two independent implementations agree: the native enumerator
visits $22\,641\,594$ nodes, while the optimizer-free Python replay
reconstructs the same complete set.  The low-word tests through degree six,
including nonnegativity of the restricted low-product remainders, leave
exactly five profiles for each character labelling, of ranks
\[
 24,\quad 30,\quad 32,\quad 35,\quad 38.
\]
These are finite necessary induction data, not fusion categories or modular
data.  In particular, the ring remains unresolved.

The height-six cyclotomic system, strengthened by
Proposition~\ref{prop:odd-column-orbits}, eliminates all five profiles for
one labelling and the rank-$24$ and rank-$32$ profiles for the other.
Each elimination is witnessed by an exact integral Farkas functional.
Thus the retained induction-profile ranks are $30$, $35$, and $38$, with
one character labelling.  The rank-$35$ and rank-$38$ systems have explicit
nonnegative integral witnesses.  A stronger rank-$35$ witness also satisfies
the low-Verlinde and restriction-product constraints.  These witnesses
verify only the stated necessary equations; they do not give the missing
center fusion coefficients or a modular $S$-matrix.
